\documentclass[10pt,a4paper]{article}
\usepackage[utf8]{inputenc}
\usepackage[a4paper,left=3cm,right=2cm]{geometry}
\usepackage{amsmath}
\usepackage[thmmarks, amsmath, thref]{ntheorem}
\usepackage{amsfonts}
\usepackage{amssymb}
\usepackage{fancyhdr}
\pagestyle{fancy}
\usepackage{anysize}
\usepackage{graphicx}
\usepackage{float}
\usepackage{hyperref}
\usepackage{enumitem}
\newtheorem{theorem}{Theorem}
\theoremsymbol{\ensuremath{\blacksquare}}
\newtheorem*{solution}{Solución:}
\usepackage{float}
\usepackage{xcolor}

% Margenes y formalidades

\textwidth 6.25in % Margen Derecho y ancho del texto
\textheight 8.5in
\oddsidemargin 0in % Margen Izquierdo
\headheight 0in % Largo del texto
\leftmargin 1in
\parindent 0em % Salto entre la indentación (Sangría)https://www.overleaf.com/4519127669sdpptkhvsbfn
\topmargin -0.5in % Margen Superior
\parskip 0.5ex % Salto entre parrafos
\baselineskip 1pt



\lhead{
\includegraphics[scale=0.09]{Figuras/logoUSM.png}} % texto izquierda de la cabecera
%\chead{}                                      % texto centro de la cabecera
\rhead{Universidad Técnica Federico Santa María\\Ingeniería Comercial } % número de página a la derecha
\renewcommand{\headrulewidth}{0.4pt} % grosor de la línea de la cabecera
\renewcommand{\footrulewidth}{0.4pt} % grosor de la línea del pie

\begin{document}
\vspace{5cm}



\begin{center}
\begin{large}
\textbf{\\CERTAMEN RECUPERATIVO}
\end{large}\\
\end{center}

\begin{center}
\begin{large}
\textbf{Econometría ICS-294}
\end{large}\\
\vspace{0.2cm}
II trimestre 2025\\
Profesor: Pedro Comte Hidalgo
\end{center}

\textbf{Instrucciones:} 

\begin{itemize}

    \item Usted dispone de 100 minutos  para intentar los 100 puntos del certamen.
    \item Cuide la presentación y redacción de sus preguntas.
    \item No puede utilizar celular, computador o tablet, pero sí una calculadora TI.
     
\end{itemize}



\section{Preguntas Cortas (40 Puntos)}

\begin{enumerate}

\item (8 Puntos) ¿Que es la heterocedasticidad y por qué en presencia de ella la inferencia estadística deja de ser consistente?

\item (8 Puntos) Suponga un modelo que mide la probabilidad de cometer un asesinato dado que la persona realizó el servicio militar. Sabiendo que existe endogeneidad, ¿cuál podría ser una variable instrumental válida para este problema y por qué? 

\item (8 Puntos) ¿Cuál es la diferencia entre multicolinealidad perfecta y alta multicolinealidad, y como afectan al calculo de los regresores mediante MCO?

\item (8 puntos) Explique brevemente por qué una variable que sigue una caminata aleatoria (\textit{random walk}) no es estacionaria.

\item (8 puntos) ¿Qué efecto tiene aplicar primeras diferencias a una serie con raíz unitaria?

\end{enumerate}

\section{Ejercicios (60 Puntos)}

\begin{enumerate}


\item (20 puntos) Se tienen los siguientes datos de años de eduación universitaria y experiencia laboral y salario de 5 individuos:
\begin{center}
\begin{tabular}{|c|c|c|c|}
\hline
\textbf{Individuo} & \textbf{Educ} ($X_1$) & \textbf{Exp} ($X_2$)  & \textbf{Salario ($Y$)}\\
\hline
1 & 0 &  8 & 1800 \\
2 & 2 &  6 & 2300 \\
3 & 4 &  5 & 2400 \\
4 & 6 &  4 & 2500 \\
5 & 8 &  3 & 3000 \\
\hline
\end{tabular}
\end{center}

\begin{enumerate}
    \item (5 puntos) Planteando un modelo donde el salario sólo dependa de la educación universitaria de la forma:
    $$Salario = \beta_0 + \beta_1 \cdot Educ$$
    Encuentre los valores de los estimadores $\hat{\beta_0}$ y $\hat{\beta_1}$ que predicen el valor del salario, e interprete su resultado.
    
    \item (5 puntos) Calcule Bondad de ajuste de este modelo, concluya.
        
    \item (5 puntos) Luego se realiza un modelo incluyendo la variable Experiencia, por lo que el salario estaría explicado por la educación universitaria y la experiencia de la siguiente manera:
    $$Salario = \beta_0 + \beta_1 \cdot Educ + \beta_2 \cdot Exp$$
    Calcule los estimadores $\hat{\beta_0}$, $\hat{\beta_1}$ y $\hat{\beta_2}$, que nos sirven para predecir el salario, sabiendo que:
    \[
(X^tX)^{-1} \;=\;
\begin{pmatrix}
145 & -11.5 & -19 \\
-11.5 & 0.925 & 1.5 \\
-19 & 1.5 & 2.5
\end{pmatrix}
\]

\item (5 puntos) Calcule Bondad de ajuste de este último modelo, concluya.
\end{enumerate}

\item (20 Puntos) Se dispone de una muestra de 250 observaciones con la cual se estimó el siguiente modelo de regresión lineal múltiple:
$$Y_i = \beta_0 + \beta_1\,\textit{L}_i + \beta_2\,\textit{K}_i + u_i$$
Donde . $Y$ es la producción, $L$ es cantidad de dinero puesta en Trabajo, y $K$ la cantidad de dinera puesta en capital.
Se desea contrastar el supuesto de heterocedasticidad y, para ello, se aplicará el método de Breusch-Pagan.
\begin{enumerate}
    \item (5 Puntos) Plantee la regresión auxiliar para el método de Breusch-Pagan\\
    \textbf{Solución}\\
    $$\hat{u}_i^{\,2} = \alpha_0 + \alpha_1 {L}_i + \alpha_2 {K}_i + v_i$$
    
    \item (5 puntos) Formule la hipótesis nula y la hipótesis alternativa para contrastar la presencia de heterocedasticidad.
    
    \item (10 Puntos) Si el $R^2$ de la regresión auxiliar es $0.05$, contraste, al 5\% de significación, la existencia de heterocedasticidad.
    
    \end{enumerate}

\item Se realiza un modelo de regresión lineal múltiple para una muestra de 209 datos que entrega el siguiente resultado:
\begin{center}
    \begin{figure}[H]
        \centering
        \includegraphics[width=0.8\linewidth]{certamenes/Certamenes T2/salida_modelo_cr.png}
    \end{figure}
\end{center}
Realice los siguientes test de hipotesis:
\begin{enumerate}
    \item La variable ROE es igual 15, con un nivel de significancia del 5\%.
    
    \item La variable ROS es menor a 0, con una significancia del 1\%.
    
    \item Que el modelo el significativo globalmente al 5\%.
    
    \end{enumerate}


\end{enumerate}


\section*{Formulario}


\begin{itemize}

\item El intervalo de confianza de $100(1-\alpha)\%$ para la media poblacional $\mu$ cuando $\sigma$ es \emph{\textbf{conocido}}:
$\color{red}{\bar{x} \pm z_{\frac{\alpha}{2}}\frac{\sigma}{\sqrt{n}}}$

\item El intervalo de confianza de $100(1-\alpha)\%$ para la diferencia de medias
$\mu_x-\mu_y$ es:
${\color{red}(\bar{X}-\bar{Y}) \pm z_{\frac{\alpha}{2}}\sqrt{\frac{\sigma_x^2}{n_x}+\frac{\sigma_y^2}{n_y}}}$

\item  Estimadores MCO.
{\color{red}
$$\hat{\beta}_0 = \bar{y} - \hat{\beta}_1 \bar{x} \quad \quad   \hat{\beta}_1 = \frac{\sum_{i=1}^{n}(x_i - \bar{x})(y_i-\bar{y})}{\sum_{i=1}^{n}(x_i - \bar{x})^2}$$}
    
\item 
	{\color{red}	\[ SST = \sum_{i=1}^{n} (y_i - \bar{y})^2 \quad \quad  SSE = \sum_{i=1}^{n} (\hat{y}_i - \bar{y})^2 \quad \quad SSR = \sum_{i=1}^{n} \hat{u}_i^2  \]}
		
		
    \item Vector de estimadores:    
{\color{red}    $\hat{\beta}=(X^TX)^{-1}X^Ty$}

\item {\color{red}\[
A = \begin{bmatrix}
a & b \\
c & d
\end{bmatrix}
\qquad
A^{-1} = \frac{1}{det(A)} \begin{bmatrix}
d & -b \\
-c & a
\end{bmatrix}
\]}
\end{itemize}

\begin{itemize}

\item Estadistico t para $\beta_j = a_j$

$$\color{red}{t_{\beta_j}=\frac{\hat{\beta_j}-a_j}{s.e.(\hat{\beta_j})}}$$

\item El intervalo de confianza de $100(1-\alpha)\%$ para el estimador $\beta_j$:

$${\color{red}\hat{\beta_j} \pm t_{\frac{\alpha}{2},n-k-1}s.e.(\hat{\beta_j})}$$


\item para el estadistico F entre dos modelos:

$$\color{red}{F=\frac{SSR_r-SSR_{ur}}{SSR_{ur}}\frac{n-k-1}{j}}$$
$$\color{red}{F=\frac{R^2_{ur}-R^2_{r}}{1-R^2_{ur}}\frac{n-k-1}{j}}$$

\item Test de Chow para Grupos.

    \[\color{red}
    F = \frac{(SSR_r - SSR_{ur}) / ((G - 1)(k + 1))}{SSR_{ur} / (n - G(k + 1))} \sim F((G - 1)(k + 1), n - G(k + 1))
    \]

Con

$$\color{red}SSR_{ur} = \sum_{g=1}^G SSR_g$$

\item Estadistico t para igualdad entre parametros:
 \[\color{red}t_{\text{estadístico}} = \frac{\hat{\beta}_1 - \hat{\beta}_2 - 0}{\sqrt{\hat{Var}(\hat{\beta}_1 - \hat{\beta}_2)}}\]


\item  Varianza de una suma.
{\color{red}
$$\text{Var}(kA + cB) = k^2\text{Var}(A) + c^2\text{Var}(B) + 2kc\text{COV}(A, B)$$}

\item $t_{0.95,n} = 1.64$ cuando $n >50$.
\item $t_{0.975,n} = 1.96$ cuando $n >50$.
\item $t_{0.99,n}=2.33$ cuando $n>50$


\item tabla Fisher para 5\%:
\begin{table}[H]
\centering
\caption{Valores críticos de la distribución F (cola superior) para $\alpha=0.05$. 
El numerador usa $df_1$ y el denominador $df_2$.}
\label{tab:Fcrit_5pct}
\begin{tabular}{|l|r|r|r|r|r|}
\toprule
\hline
& \multicolumn{5}{c|}{$df_2$} \\
\hline
$df_1$ & 50 & 100 & 150 & 200 & 250 \\
\hline
\midrule
1 & 4.0343 & 3.9361 & 3.9042 & 3.8884 & 3.8789 \\
\hline
2 & 3.1826 & 3.0873 & 3.0564 & 3.0411 & 3.0319 \\
\hline
3 & 2.7900 & 2.6955 & 2.6649 & 2.6498 & 2.6407 \\
\hline
4 & 2.5572 & 2.4626 & 2.4320 & 2.4168 & 2.4078 \\
\hline
5 & 2.4004 & 2.3053 & 2.2745 & 2.2592 & 2.2501 \\
\end{tabular}
\end{table}

\begin{itemize}

\item Tabla $\chi^2$
\begin{center}
    \begin{tabular}{|c|c|c|c|}
    \hline
        df & 0.99 & 0.95 & 0.90 \\ 
        \hline
        1 &  6.635 & 3.841 & 2.706\\
        \hline
        2 &  9.210 & 5.991 & 4.605 \\
        \hline
        3 & 11.345 & 7.815 & 6.251  \\
        \hline
        4 &  13.277 & 9.488 & 7.779 \\
        \hline
        5 &  15.086 & 11.070 & 9.236 \\
        \hline
        6 &  16.812 & 12.592 & 10645 \\
        \hline
    \end{tabular}
\end{center}

\item Multiplicado de Lagrange

$$\color{red}{LM = nR^2 \sim \chi^2}$$

\item Test F para regresión auxiliar.
\begin{center}
 {  \color{red}     \(F=\dfrac{R_{\text{aux}}^{2}/q}{(1-R_{\text{aux}}^{2})/(n-k-1)}
          \sim F_{q,\;n-k-1}\).}
\end{center}



\item Tabla de test D-F
\begin{center}
    \begin{figure}[H]
        \centering
        \includegraphics[width=0.7\linewidth]{sem14/tabla.png}
    \end{figure}
\end{center}

\end{itemize}

\end{itemize}
\end{document}